\begin{definition}[The graded-agreement specification]
  \label{def:aba-gbca-spec}
  \lean{PLTS.ABA.GBCA.Step}\lean{PLTS.ABA.GBCA.specFamily}\leanok
  \uses{def:aba-labels, def:idle-family}
  Transition System 2, as a round-indexed instance and the $\mathbb{N}$-indexed family of Definition~\ref{def:idle-family} over it: internal binding, the graded returns $(b,2)$, $(b,1)$ and $(\bot,0)$, and the announcement of the bound bit. The family is an LTS.\lean{PLTS.ABA.GBCA.specFamily_isLTS}\leanok{} An instance binds by exclusion (D19): its state carries the set $excluded \subseteq \{0,1\}$ of bits the instance can no longer return, written only by an internal exclude enabled from $excluded = \emptyset$ alone, so the exclusion fires at most once per instance; beside it the state carries the pending call $call$ and the return flag $ret$ of each process, the grade $grade$, set by a grade-2 return and by a grade-0 return exclusively, and the corrupted set.
  That exclude, the dissenting bit at a grade-1 return and each of the two bits at a grade-0 return carry an $f{+}1$ $F$-blind support guard, $\mathrm{InputSupport}$: $f{+}1$ distinct callers-or-corrupted of the bit in question (D14, D15).
  The returns of grade 1 and 2 carry the guard pair $v \notin excluded$ and $1-v \in excluded$, and the grade-0 return carries $excluded \neq \emptyset$, so every return requires an excluded bit; each return label $\mathrm{retG}(r,id,out)$ carries the graded outcome alone.
  The announcement \leandeclnamed{announceBound}{PLTS.ABA.GBCA.Step.announceBound} fires on $\mathrm{boundG}(r,b)$ under $1-b \in excluded$, writes nothing and may repeat, so it is enabled at every state from the round's first return on.
  On a reachable state $excluded$ holds at most one bit, so that guard determines $b$ as the complement of the excluded bit and all the announcements of a round name one bit, and a return of $v$ at grade 1 or 2 returns the announced bit.
  The bit is a ghost output (D29): it is the bound value the specification holds as state, and no program reads it.
  In a round that returns no value the run ends with grade-0 returns alone, and the announcement names the bit not excluded, so the round is bound on the trace all the same.
  The call label is taken at every state, by \leandeclnamed{call}{PLTS.ABA.GBCA.Step.call} at a process holding no call and by \leandeclnamed{callLoop}{PLTS.ABA.GBCA.Step.callLoop} at every state, and $\mathrm{fail}$ by a total transition.
  The full list of transitions: \leandecl{PLTS.ABA.GBCA.specFamily}.
\end{definition}
